% ====================================================================
% Файл: tasks/01_mechanics/kim03.tex
% Тема: Задание №3 (Законы сохранения в механике)
% ====================================================================

% === ЗАДАЧА 1 ===
% Тег: #МЕХАНИКА #КИМ_03 #ВАРИАНТ_1
\begin{taskbasic}[code={\label{task:dem26_v1_n3}}]
\textbf{Задача 1.} При равномерном прямолинейном перемещении саней по горизонтальному участку пути на $20$~м постоянная горизонтально направленная сила тяги совершает работу $240$~Дж. Чему равен модуль силы трения?

\kim{3}
\end{taskbasic}
\answer{12~Н}

% === ЗАДАЧА 2 ===
% Тег: #МЕХАНИКА #КИМ_03 #ВАРИАНТ_2
\begin{taskbasic}[code={\label{task:dem26_v2_n3}}]
\textbf{Задача 2.} Под действием постоянной горизонтально направленной силы тяги $80$~Н ящик равномерно прямолинейно перемещается по горизонтальному участку пути на расстояние $15$~м. Какую работу совершает при этом сила трения?

\kim{3}
\end{taskbasic}
\answer{-1200~Дж}

% === ЗАДАЧА 3 ===
% Тег: #МЕХАНИКА #КИМ_03 #ВАРИАНТ_3
\begin{taskbasic}[code={\label{task:dem26_v3_n3}}]
\textbf{Задача 3.} Тело массой $10$~кг движется в инерциальной системе отсчёта по прямой. Под действием постоянной силы, направленной вдоль этой прямой, за $4$~с импульс тела увеличился с $12$~\text{кг}$\cdot$\text{м/с} до $48$~\text{кг}$\cdot$\text{м/с}. Определите величину этой силы.

\kim{3}
\end{taskbasic}
\answer{9~Н}

% === ЗАДАЧА 4 ===
% Тег: #МЕХАНИКА #КИМ_03 #ВАРИАНТ_4
\begin{taskbasic}[code={\label{task:dem26_v4_n3}}]
\textbf{Задача 4.} Тело массой $10$~кг движется в инерциальной системе отсчёта по прямой. Под действием постоянной силы величиной $9$~Н, направленной вдоль этой прямой, за $2$~с импульс тела увеличился и стал равен $24$~\text{кг}$\cdot$\text{м/с}. Определите начальный импульс тела.

\kim{3}
\end{taskbasic}
\answer{6~\text{кг}$\cdot$\text{м/с}}

% === ЗАДАЧА 5 ===
% Тег: #МЕХАНИКА #КИМ_03 #ВАРИАНТ_5
\begin{taskbasic}[code={\label{task:dem26_v5_n3}}]
\textbf{Задача 5.} Скорость волана при подаче в бадминтоне составила $288$~км/ч. Масса волана $6$~г. Какова кинетическая энергия волана при подаче?

\kim{3}
\end{taskbasic}
\answer{24~Дж}

% === ЗАДАЧА 6 ===
% Тег: #МЕХАНИКА #КИМ_03 #ВАРИАНТ_6
\begin{taskbasic}[code={\label{task:dem26_v6_n3}}]
\textbf{Задача 6.} Какую скорость при подаче сообщили волейбольному мячу, если его кинетическая энергия в этот момент составила $56$~Дж. Масса мяча $280$~г.

\kim{3}
\end{taskbasic}
\answer{20~м/с}

% === ЗАДАЧА 7 ===
% Тег: #МЕХАНИКА #КИМ_03 #ВАРИАНТ_7
\begin{taskbasic}[code={\label{task:dem26_v7_n3}}]
\textbf{Задача 7.} У основания гладкой наклонной плоскости шайба обладает скоростью, модуль которой равен $4$~м/с. Определите массу шайбы, если максимальная потенциальная энергия, которую она приобретает при подъёме по плоскости относительно её основания, составляет $0{,}4$~Дж. Сопротивлением воздуха пренебречь.

\kim{3}
\end{taskbasic}
\answer{0{,}05~кг}

% === ЗАДАЧА 8 ===
% Тег: #МЕХАНИКА #КИМ_03 #ВАРИАНТ_8
\begin{taskbasic}[code={\label{task:dem26_v8_n3}}]
\textbf{Задача 8.} У основания гладкой наклонной плоскости шайба массой $120$~г обладает скоростью, модуль которой равен $5$~м/с. Определите максимальную потенциальную энергию, которую приобретёт шайба при подъёме по плоскости относительно её основания. Сопротивлением воздуха пренебречь.

\kim{3}
\end{taskbasic}
\answer{1{,}5~Дж}

% === ЗАДАЧА 9 ===
% Тег: #МЕХАНИКА #КИМ_03 #ВАРИАНТ_9
\begin{taskbasic}[code={\label{task:dem26_v9_n3}}]
\textbf{Задача 9.} При упругой деформации, равной $2$~см, потенциальная энергия пружины равна $10$~мДж. На сколько увеличится потенциальная энергия этой пружины при увеличении упругой деформации ещё на $2$~см?

\kim{3}
\end{taskbasic}
\answer{30~мДж}

% === ЗАДАЧА 10 ===
% Тег: #МЕХАНИКА #КИМ_03 #ВАРИАНТ_10
\begin{taskbasic}[code={\label{task:dem26_v10_n3}}]
\textbf{Задача 10.} При упругой деформации, равной $1$~см, потенциальная энергия пружины равна $5$~мДж. Чему будет равна потенциальная энергия этой пружины при увеличении упругой деформации в $3$ раза?

\kim{3}
\end{taskbasic}
\answer{45~мДж}

% === ЗАДАЧА 11 ===
% Тег: #МЕХАНИКА #КИМ_03 #ВАРИАНТ_11
\begin{taskbasic}[code={\label{task:dem26_v11_n3}}]
\textbf{Задача 11.} Тело массой $200$~г, брошенное вертикально вверх с поверхности Земли, в момент броска обладало кинетической энергией, равной $20$~Дж. На какую максимальную высоту поднялось тело? Сопротивлением воздуха пренебречь.

\kim{3}
\end{taskbasic}
\answer{10~м}

% === ЗАДАЧА 12 ===
% Тег: #МЕХАНИКА #КИМ_03 #ВАРИАНТ_12
\begin{taskbasic}[code={\label{task:dem26_v12_n3}}]
\textbf{Задача 12.} Тело массой $200$~г брошено со скоростью $8$~м/с вертикально вниз с высоты $10$~м относительно поверхности Земли. Определите кинетическую энергию тела в момент падения на Землю. Сопротивлением воздуха пренебречь.

\kim{3}
\end{taskbasic}
\answer{26~Дж}

% === ЗАДАЧА 13 ===
% Тег: #МЕХАНИКА #КИМ_03 #ВАРИАНТ_13
\begin{taskbasic}[code={\label{task:dem26_v13_n3}}]
\textbf{Задача 13.} В инерциальной системе отсчёта тело массой $50$~кг движется по прямой в одном направлении под действием постоянной равнодействующей внешних сил, равной по модулю $20$~Н. Каков модуль изменения скорости тела за $8$~с?

\kim{3}
\end{taskbasic}
\answer{3{,}2~м/с}

% === ЗАДАЧА 14 ===
% Тег: #МЕХАНИКА #КИМ_03 #ВАРИАНТ_14
\begin{taskbasic}[code={\label{task:dem26_v14_n3}}]
\textbf{Задача 14.} В инерциальной системе отсчёта тело движется по прямой в одном направлении под действием постоянной равнодействующей внешних сил, равной по модулю $32$~Н. Каков модуль изменения импульса тела за $8$~с?

\kim{3}
\end{taskbasic}
\answer{256~\text{кг}$\cdot$\text{м/с}}

% === ЗАДАЧА 15 ===
% Тег: #МЕХАНИКА #КИМ_03 #ВАРИАНТ_15
\begin{taskbasic}[code={\label{task:dem26_v15_n3}}]
\textbf{Задача 15.} Тело массой $0{,}4$~кг движется в инерциальной системе отсчёта по прямой в одном направлении. При этом равнодействующая всех сил, действующих на тело, постоянна и равна по модулю $8$~Н. За какое время модуль изменения импульса тела составит $20$~\text{кг}$\cdot$\text{м/с}?

\kim{3}
\end{taskbasic}
\answer{2{,}5~с}

% === ЗАДАЧА 16 ===
% Тег: #МЕХАНИКА #КИМ_03 #ВАРИАНТ_16
\begin{taskbasic}[code={\label{task:dem26_v16_n3}}]
\textbf{Задача 16.} Тело движется в инерциальной системе отсчёта по прямой в одном направлении. При этом равнодействующая всех сил, действующих на тело, постоянна и равна по модулю $8$~Н. Каков модуль изменения импульса тела за $4$~с?

\kim{3}
\end{taskbasic}
\answer{32~\text{кг}$\cdot$\text{м/с}}

% === ЗАДАЧА 17 ===
% Тег: #МЕХАНИКА #КИМ_03 #ВАРИАНТ_17
\begin{taskbasic}[code={\label{task:dem26_v17_n3}}]
\textbf{Задача 17.} Тело движется в инерциальной системе отсчёта по прямой в одном направлении. Равнодействующая всех сил, действующих на тело, постоянна и равна по модулю $0{,}75$~Н. Модуль импульса тела изменился на $3$~\text{кг}$\cdot$\text{м/с}. Сколько времени потребовалось для этого?

\kim{3}
\end{taskbasic}
\answer{4~с}

% === ЗАДАЧА 18 ===
% Тег: #МЕХАНИКА #КИМ_03 #ВАРИАНТ_18
\begin{taskbasic}[code={\label{task:dem26_v18_n3}}]
\textbf{Задача 18.} Тело движется в инерциальной системе отсчёта по прямой в одном направлении. Равнодействующая всех сил, действующих на тело, постоянна. Модуль импульса тела за $5$~с изменился на $8$~\text{кг}$\cdot$\text{м/с}. Определите модуль равнодействующей сил, действующих на тело.

\kim{3}
\end{taskbasic}
\answer{1{,}6~Н}

% === ЗАДАЧА 19 ===
% Тег: #МЕХАНИКА #КИМ_03 #ВАРИАНТ_19
\begin{taskbasic}[code={\label{task:dem26_v19_n3}}]
\textbf{Задача 19.} Отношение модуля импульса микроавтобуса к модулю импульса мотоцикла $\frac{p_1}{p_2}=8$. Каково отношение массы микроавтобуса к массе мотоцикла $\frac{m_1}{m_2}$, если отношение модулей их скоростей $\frac{v_1}{v_2}=5$?

\kim{3}
\end{taskbasic}
\answer{1{,}6}

% === ЗАДАЧА 20 ===
% Тег: #МЕХАНИКА #КИМ_03 #ВАРИАНТ_20
\begin{taskbasic}[code={\label{task:dem26_v20_n3}}]
\textbf{Задача 20.} Отношение модуля импульса легкового автомобиля к модулю импульса мотоцикла $\frac{p_1}{p_2}=6$. Каково отношение модулей их скоростей $\frac{v_1}{v_2}$, если отношение массы легкового автомобиля к массе мотоцикла $\frac{m_1}{m_2}=2$?

\kim{3}
\end{taskbasic}
\answer{3}